% Written for Swift Shifter's benchmark (#150); part of the repository,
% under its licence (GPL-3.0-or-later).
\documentclass[10pt,twocolumn]{ctexart}
\usepackage{amsmath}
\usepackage{geometry}
\geometry{margin=2cm}

\title{雙欄排版下的閱讀順序判斷}
\author{Swift Shifter 測試文件}
\date{2026 年 10 月}

\begin{document}
\maketitle

\begin{abstract}
從 PDF 擷取文字時，雙欄頁面的閱讀順序是最常見的錯誤來源。本文描述一個以文字座標為基礎的
簡單方法：先找出欄間的空白帶，再依欄位由上而下排序。我們也討論 header、footer 與跨欄標題
等例外情況。
\end{abstract}

\section{問題}

PDF 只記錄每個字形在頁面上的位置，並不記錄段落或閱讀順序。對單欄頁面，依 $y$ 座標排序
通常就夠了；但在雙欄頁面上，同一高度的兩行分屬不同欄位，若只依 $y$ 排序，左右兩欄的文字
會被交錯合併，產生難以閱讀的結果。

這個問題在學術論文中特別常見。多數會議論文採用雙欄格式，而且常有跨欄的標題、圖表與公式。

\section{方法}

\subsection{找出欄間空白}

把頁面沿 $x$ 軸切成寬度為 $w$ 的長條，計算每一條中被文字覆蓋的比例 $c(x)$。欄間空白是
頁面中段 $c(x)$ 接近零的連續區間：
\begin{equation}
  G = \{\, x \mid c(x) < \epsilon,\ x_{\min} < x < x_{\max} \,\}.
\end{equation}
若 $G$ 為空，就把頁面當成單欄。

\subsection{排序}

找到空白帶的中心 $g$ 之後，每個文字區塊依其中心點 $x_b$ 分到左欄或右欄，再分別依 $y$
排序。跨過 $g$ 的區塊，例如標題或寬表格，則依其 $y$ 座標插在兩欄之間，並把頁面在該處切成
上下兩段，各自重複上述步驟。

\subsection{Header 與 footer}

出現在每頁相同位置、而且內容相似（例如頁碼或期刊名稱）的區塊，視為 header 或 footer，
不放進正文。判斷時比較連續數頁中位置差在 $\delta$ 以內的區塊，若文字的 edit distance 小於
門檻就移除。

\section{評估}

我們以原始碼為標準答案，計算擷取結果的 normalized edit distance：
\begin{equation}
  d(a, b) = \frac{\operatorname{lev}(a, b)}{\max(|a|, |b|)}.
\end{equation}
分數 $1 - d$ 越接近 1 越好。初步結果顯示，加入欄位判斷後，雙欄頁面的分數明顯提高，而單欄
頁面幾乎不受影響。

\section{限制}

這個方法假設欄位是規則的直條，對不規則的雜誌排版、文繞圖，或一頁中欄數改變的情況效果較差。
這些情況需要更完整的版面分析，例如以機器學習模型辨識區塊類型。

\begin{thebibliography}{9}
\bibitem{lev} V. I. Levenshtein. Binary codes capable of correcting deletions,
  insertions, and reversals. \emph{Soviet Physics Doklady}, 10(8):707--710, 1966.
\bibitem{breuel} T. M. Breuel. Two geometric algorithms for layout analysis.
  In \emph{Document Analysis Systems V}, pages 188--199, 2002.
\end{thebibliography}

\end{document}
